\begin{DndMonster}[width=0.5\textwidth]{Blood Leech\phantomsection\addcontentsline{toc}{subsection}{Blood Leech}\label{monster:BloodLeech}}
	\DndMonsterType{Tiny Monstrosity, Unaligned}
	
	% If you want to use commas in the key values, enclose the values in braces.
	\DndMonsterBasics[
		armor-class = {10},
		initiative	= +0,
		hit-points  = {\DndDice{8d4}},
		speed       = {20 ft., Climb 20 ft., Swim 30 ft.},
	]
	
	\renewcommand{\AbilityScoreSpacer}{~}
	\DndMonsterAbilityScores[
		str = 13,
		dex = 10,
		con = 10,
		int = 10,
		wis = 12,
		cha = 5,
	]
	
	\DndMonsterDetails[
%		saving-throws			= {CON +3, WIS +2},
%		skills					= {Perception +6},
%		damage-vulnerabilities	= {Cold},
%		damage-resistances		= {Acid, Bludgeoning, Cold, Fire, Lightning, Piercing, Slashing, Thunder},
%		damage-immunities		= {Necrotic, Poison},
%		gear					= {Leather Armor, Light Crossbow, Scimitar},
		senses					= {Blindsight 30 ft., Tremorsense 60 ft., Passive Perception 11},
%		condition-immunities	= {Charmed, Exhaustion, Frightened, Grappled, Paralyzed, Petrified, Poisoned, Prone, Restrained},
%		languages				= {-},
		challenge				= 1,
		proficiency-bonus		= +2,
	]
	
	\DndMonsterSection{Traits}
	\DndMonsterAction{Aquatic Camouflage}
	The leech has Advantage on Dexterity (Stealth) checks made to hide in aquatic terrain.
	
	\DndMonsterAction{Blood Bags}
	The leech has three blood bags in its body. When encountered, roll 1d3 to determine how many blood bags are filled with blood. The rest are empty.
	
	\DndMonsterSection{Actions}
	\DndMonsterAttack[
		name			= Attach,
		distance		= melee,	% valid options are in the set {both,melee,ranged}
%		type			= weapon,	% valid options are in the set {weapon,spell}
		mod				= +3,
		reach			= 0,
%		range			= 20/60,
		targets			= one target,
%		dmg				= {},
%		dmg-type		= ,
%		plus-dmg		= \DndDice{1d6},
%		plus-dmg-type	= Poison,
%		or-dmg			= \DndDice{1d20},
%		or-dmg-when		= {},
		extra			= {The leech attaches to the target. While attached, the leech doesn't attack, and moves with the target, using none of the leech's movement. The leech can detach itself using 5 feet of movement, and detaches immediately if its blood bags are filled or the target dies.\\A creature within 5 feet can take an action to detach the leech with a successful DC 13 Wisdom (Medicine) check.\\Failure: The leech is still detached, but the target has the Bleeding condition (2d4).\\Failure by 5 or More: The leech isn't detached and the target takes \DndDice{2d4 + 3} Piercing damage},
	]
	
	\DndMonsterSection{Bonus Actions}
	\DndMonsterAction{Blood Drain}
	The leech can use this Bonus Action only when it is attached to a creature. The attached creature loses \DndDice{3d6} Hit Points due to blood loss, and the blood leech fills one of its empty blood bags with some of the target's blood.
	
	\DndMonsterAction{Blood Transfusion}
	The leech can use this Bonus Action only when it is attached to a creature and has at least one filled blood bag. The leech regains \DndDice{3d6} Hit Points, and it empties one of its filled blood bags.
\end{DndMonster}